\section{Parte I. Preparación de los datos}
\label{sec:parte_01}

\subsection{Elección de variables}
A partir de la base de datos de la Encuesta Permanente de Hogares (EPH) correspondientes al primer trimestre de 2026, se selecionaron únicamente a los ocupado (\texttt{ESTADO=1}) a quienes luego se los clasificó como formales (\texttt{EMPLEO=1}) e informales (\texttt{EMPLEO=2}). Entonces, para la construcción de la variable predicha, únicamente se recodificó la variable \texttt{EMPLEO} de la siguiente forma:
\begin{equation}
\mathrm{informal}_i =
\left\{\begin{array}{ll}
1 & \mbox{si } \mathrm{EMPLEO}_i=2,\\
0 & \mbox{si } \mathrm{EMPLEO}_i=1.
\end{array}\right.
\label{eq:informal_parte01}
\end{equation}
El conteo inicial encontró que existen 19.560 trabajadores, de los cuales el 55,02\% (10.762) corresponde a formales y el 44,98\% (8.798) a informales. 

La selección de las variables (\tabref{tab:variables_parte01}) se realizó a partir de una breve revisión de literatura apoyado por la IA, aunque las variables de residencia hace 5 años (CH16) y régimen de tenencia de la vivienda (II7) fueron directamente selecionadas por la IA. En general estas variables contemplan características demográficas, educativas, laborales, territoriales y del hogar que permitirán predecir la informalidad basada en trabajos previos.
\begin{table}[H]
	\centering
	\caption{Set de variables seleccionadas}
	\label{tab:variables_parte01}
	\footnotesize
	\begin{tabular}{p{1.8cm}p{5cm}p{7cm}}
		\hline
		\textbf{Variable} & \textbf{Característica} & \textbf{Antecedente} \\
		\hline
		\texttt{CH03} & Parentesco con la persona de referencia del hogar & \citet{casal2013} \\
		\texttt{CH04} & Sexo & \multirow{2}{7cm}{\citet{formichella2008,ludmer2020}} \\
		\texttt{CH06} & Edad & \\
		\texttt{CH07} & Situación conyugal & \citet{casal2013} \\
		\texttt{CH15} & Lugar de nacimiento & \citet{ludmer2020} \\
		\texttt{CH16} & Lugar de residencia cinco años antes & Propuesta por la IA \\
		\texttt{NIVEL\_ED} & Nivel educativo alcanzado & \citet{formichella2008,alvarez2025} \\
		\texttt{INTENSI} & Intensidad de la ocupación & \citet{alvarez2025} \\
		\texttt{PP04A} & Tipo de establecimiento & \citet{bertranou2013} \\
		\texttt{PP04C} & Tamaño del establecimiento & \citet{cruces2018,ludmer2020} \\
		\texttt{P21} & Ingreso de la ocupación principal & \citet{alvarez2025} \\
		\texttt{REGION} & Región estadística & \multirow{2}{7cm}{\citet{cristina2019}} \\
		\texttt{MAS\_500} & Tamaño del aglomerado & \\
		\texttt{IX\_MEN10} & Menores de 10 años en el hogar & \citet{casal2013} \\
		\texttt{II7} & Régimen de tenencia de la vivienda & Propuesta por la IA \\
		\hline
	\end{tabular}
\end{table}
\vspace{-0.5cm}
La consigna indicaba que se debía evitar la elección de las variables que predecían el \texttt{EMPLEO}, ya que el modelo utilizaría información que, prácticamente revelaría la clasificación que se busca predecir; es decir, se generaría un poder predictivo artificialmente elevado.

\subsection{Construcción de variables}
Sumado a lo anterior, se construyeron tres variables adicionales basadas en la propia recomendación de la consigna y estos fueron edad al cuadrado, tamaño del hogar y la proporción de integrantes ocupados. Adicionalmente, se incluyó el logaritmo del ingreso y un indicador de ingreso no declarado para distinguir los valores faltantes de los valores 0 cero.

\subsection{Análisis de valores faltantes}
Algunas variables presentaban un valor diferente que; por lo general, estaban representados por 9, 99 o $-9$, en cuyos casos se reeamplazaron por \texttt{NaN} segun corresponda. En la \figref{fig:heatmap_faltantes} se presentan los valores \texttt{NaN} por variable y se observa que el ingreso (P21) y en el número de personas que trabajan en el centro laboral (PP04C) concentran la mayor cantidad de \texttt{NaN}, que fueron 3719 y 2607 respectivamente.
\begin{figure}[H]
	\centering
	\caption{Heatmap de valores faltantes}
	\includegraphics[width=0.91\textwidth]{\joinpath{\outputspath}{B_Figure/heatmap_faltantes.png}}
	\label{fig:heatmap_faltantes}
\end{figure}
\vspace{-0.5cm}
Mediante la prueba $t$ de Welch que es utilizada por \texttt{scipy.stats.ttest_ind()}, se comparan las medias de los trabajadores informales con las de los formales y algunos resultados se presentan en la \tabref{tab:dif_medias_resumen}.
\input{\joinpath{\outputspath}{B_Table/tabla_diferencia_medias_resumen.tex}}
\vspace{-0.5cm}
En general, se puede concluir que existen diferencias estadísticamente signifivaticas entre los trabajadores formales e informales, aunque estos resultados no deben interpretarse como relaciones causales. Por ejemplo, trabajadores informales son, en promedio, 4 años menores a los formales o que, en promedio, residen en hogares más numerosos y registran menores ingresos laborales. No obstante, la diferencia entre las proporciones de mujeres no resulta estadísticamente significativa.

\subsection{Train y Test split}
La muestra final estuvo compuesta por 19.560 trabajadores (observaciones), de los cuales 13.692 corresponden al train group (70\%) y 5.868 al test group (30\%). Esta división permitirá evaluar posteriormente la capacidad de los modelos para predecir la informalidad laboral en observaciones que no fueron utilizadas durante el entrenamiento.